\section{Discussion}
\label{sec:discussion}

\subsection{Composite Scoring Versus Single-Metric Ranking}
The most consequential empirical observation in Section~\ref{sec:results} is that XGBoost achieved the lowest cross-validated RMSE (1109.45) among all twelve successfully trained models, yet ranked ninth of twelve under the composite score of Eq.~\eqref{eq:composite-score}, while $k$-nearest-neighbor regression, with a marginally higher RMSE (1162.41), was selected as the overall best model. This divergence is a direct and intended consequence of the composite score weighting mean absolute error, mean absolute percentage error, forecast stability, and computation time alongside RMSE, rather than ranking on RMSE alone. Inspecting Table~\ref{tab:model-comparison}, XGBoost's per-fold training time (1.97~s) is roughly three orders of magnitude larger than KNN's (0.0001~s), and its relatively higher forecast-stability penalty on this particular volatile series contributed further to its lower composite ranking despite its RMSE advantage. Whether this outcome is desirable depends on the deployment context: an operator who cares only about point-forecast accuracy and has no computation-time constraint might reasonably prefer XGBoost's result, whereas the system's default composite criterion favors a more holistic notion of a "good" model that also considers stability and cost, consistent with the broader argument in the forecasting-evaluation literature that a single error metric is an incomplete criterion for practical model selection~\cite{petropoulos2022forecasting,januschowski2020criteria}. Because both the full leaderboard and every component of the composite score are retained and exposed, a user who disagrees with the default weighting can still identify and act on the RMSE-optimal alternative directly from the returned leaderboard.

\subsection{Interpreting the Negative $R^2$ Results}
Ten of the twelve models produced a negative cross-validated $R^2$ on the evaluation dataset (Section~\ref{sec:results}), indicating that, within the walk-forward cross-validation folds, most candidate models performed worse in squared-error terms than a naive constant forecast at the training-window mean. This is best interpreted as a property of the particular evaluation dataset rather than of the system: the underlying transaction data, aggregated from 1{,}000 individually recorded sales with no explicit trend or seasonal design, exhibits substantial day-to-day volatility relative to any exploitable autocorrelation structure, a pattern common in real, unaugmented retail transaction logs that lack accompanying promotional, price, or holiday context. The honest reporting of this outcome, rather than selectively reporting only favorable results, is consistent with the trustworthiness objective stated in Section~\ref{sec:problem}, and the same evaluation would be expected to yield substantially higher $R^2$ values on a dataset with genuine, learnable trend or seasonal structure, as informally observed during system development on synthetic datasets constructed with explicit trend and promotional signal.

\subsection{Advantages of the Proposed Approach}
The breadth of the model ensemble is the system's principal practical advantage: because model families differ systematically in which data characteristics they exploit well, as documented extensively in the forecasting-competition literature~\cite{makridakis2020m4,makridakis2022m5}, training all nineteen candidates and adjudicating automatically removes the need for a user, who may have no data-science background, to correctly anticipate in advance which single technique suits their specific dataset. The frequency-aware deep-learning gate further prevents a subtle failure mode in which an under-trained neural network appears in the leaderboard alongside adequately-fitted classical models with no indication that its apparent performance is unreliable. Finally, deriving forecast uncertainty from each model's own cross-validation error, rather than a fixed percentage band applied uniformly, ties the reported confidence directly to how well that specific model actually fit that specific dataset.

\subsection{Limitations}
Several limitations follow directly from the results and the design choices documented in this paper. First, the gradient boosting, random forest, and deep learning architectures use fixed hyperparameters rather than per-dataset tuning (Section~\ref{sec:experimental}), and per-dataset tuning could plausibly improve their reported accuracy at additional computational cost. Second, MAPE, one of the three error components in the composite score, is known to be unstable near zero-valued observations, as directly observed in Section~\ref{sec:results}; a scale-free alternative such as a mean absolute scaled error would be less sensitive to this effect and is a natural refinement. Third, exogenous feature support is currently limited to a single price column and a single binary promotional flag, identified by column-name matching, rather than arbitrary user-specified exogenous variables. Fourth, the system does not currently reconcile forecasts hierarchically across a store-region-total or product-category structure, so per-group forecasts obtained through the grouping mechanism are not guaranteed to sum consistently to an aggregate forecast obtained separately.

\subsection{Scalability and Deployment Feasibility}
The concurrent thread-pool orchestration and per-model timeout mechanism allow the system to bound total response latency regardless of which subset of models eventually succeeds, and the frequency-aware gating and adaptive cross-validation fold count similarly bound cost as a function of dataset size rather than applying a fixed, potentially excessive evaluation budget uniformly. The file-based storage layer (Section~\ref{sec:proposed}) is adequate for a single-process deployment but is an explicit scalability limitation: the in-memory forecast cache and session state are local to one backend process and are not shared across multiple horizontally-scaled instances behind a load balancer, and doing so would require migrating this state to a shared store, discussed further as future work in Section~\ref{sec:future}. The largest single deployment consideration identified during this work is the size of the PyTorch dependency required for the deep learning model family, which meaningfully increases backend installation size and should be accounted for when selecting a hosting platform.

\subsection{Business Impact}
For a retail operator, the practical value of the system demonstrated in this paper is not that any single model is guaranteed superior to hand-selected alternatives, but that a non-specialist user obtains, from a single uploaded file, a fair comparison across a broad set of established techniques, an explicit and auditable rationale for which one was selected, a forecast uncertainty band grounded in that model's own demonstrated accuracy, and a natural-language narrative explicitly constrained to remain consistent with those same computed numbers. This combination directly targets the adoption barrier identified in Section~\ref{sec:problem}: a forecast that a retail manager can inspect, question through the conversational interface, and receive a consistent, grounded answer to, is more likely to be trusted and acted upon than an equally accurate forecast delivered without such support.
